\section{Cache Information: Updates and Reachable Separations}
\label{app:graph-information}
The cache graph changes the bill of an evaluation through the intersection of its demand with saved resources. Whether the full incidence relation must be retained depends on the demands that future evaluations can make. A full-panel protocol and a query-selective protocol have different sufficient statistics.


Thus full-panel trials admit an exact prefix-level implementation. Query-dependent weights and nonrectangular warm caches also admit this quotient. To determine which information selective evaluation requires, we characterize a larger class of exact compressions.



The theorem identifies the coarsest partition-based residual-sum summary that
works for every unit-price cache state. A single full panel has one group;
the three cyclic pairs $\{q_1,q_2\},\{q_2,q_3\},\{q_3,q_1\}$ have three.
Lossless bitsets, demand-profile tables, and incidence graphs implement the
same information. The partition converse quantifies over arbitrary valid warm caches on
one-step prefixes. The reachable example below establishes information loss
under the overlapping-panel protocol itself. More general encodings than group
sums can preserve the same information.
Physical restoration still requires the saved execution objects; the quotient
compresses cost accounting, not their contents.
The residual-sum formula also applies to the predictor's prices at a fixed
decision. A state-conditioned predictor refreshes those prices from saved
execution observations. Its feature state and the residual accounting table
have distinct roles; the fixed-weight publication update above concerns the latter.
For live prices $\widehat w_t$, write
$r_t(u)=\sum_q\ind\{(q,u)\notin K_t\}\widehat w_t(q,u)$.
After publishing $\mathcal Q\times P_x$ and revising prices, the exact update is
\begin{equation}
\begin{split}
 r_{t+1}(u)&=0 \quad (u\in P_x),\\
 r_{t+1}(u)&=r_t(u)+\sum_{q\in\mathcal Q}\ind\{(q,u)\notin K_t\}
       \bigl(\widehat w_{t+1}(q,u)-\widehat w_t(q,u)\bigr)
       \quad (u\notin P_x).
\end{split}
\label{eq:live-residual-update}
\end{equation}
This identity, checked in \texttt{FullPanelLivePrices.lean}, separates new
cache hits from new information about future execution. A saved stop state
can lower an uncached descendant's predicted price. Computing that revision
requires the predictor's observations. For the main fixed panel, the resulting
residual totals suffice to compute completion costs under any supplied workflow law.

\begin{proposition}[Query identity can be decision-relevant]
\label{prop:graph-separation}
There are two prefix-complete caches with the same per-prefix hit counts and the same full-panel residual bills, for which the cost-minimizing action for one fixed active query is different. No deterministic policy using only these identical summaries is optimal in both states. A randomized policy with the same action distribution in both states incurs at least $1/2$ unit of expected excess cold cost in one state. A single estimate of the same action's bill has absolute error at least $1/2$ in one state.
\end{proposition}
Take two queries $q_0,q_1$, two one-step terminal prefixes $u,v$, unit resource weights, and caches
\[
 K_L=\{(q_0,u),(q_1,v)\},\quad
 K_R=\{(q_1,u),(q_0,v)\}.
\]
Each prefix has one hit and one unit of full-panel residual cost in either state. For a demand on $q_0$, the bills of $(u,v)$ are $(0,1)$ under $K_L$ and $(1,0)$ under $K_R$. With identical candidate quality, the preferred action reverses. Adding one unit of recurring work and imposing a remaining budget of one makes only $u$ feasible in the first state and only $v$ feasible in the second. This is a finite-budget information separation, not a performance theorem for the current UCT index.

\begin{proposition}[Separation reachable from an empty cache]
\label{prop:reachable-separation}
Let the allowed panels include $A=\{q_0,q_1\}$ and $B=\{q_1,q_2\}$,
with two one-step execution prefixes $u,v$ and unit prices.
Starting empty, the complete evaluations $(A,u),(B,v)$ produce
$K_L=(A\times\{u\})\cup(B\times\{v\})$; the evaluations
$(A,v),(B,u)$ produce $K_R=(B\times\{u\})\cup(A\times\{v\})$.
Both histories use the same panel order and cost four. Both have two hits and one residual unit at each
prefix over the three-query universe. For active panel $A$, the completion
bills of $(u,v)$ are $(0,1)$ in $K_L$ and $(1,0)$ in $K_R$.
With one unit of common overhead and one unit of remaining budget, the unique
affordable action reverses.
\end{proposition}
The proof enumerates the four publications in each history and their set
differences. All five claims, including reachability and equal warm-up expenditure,
are checked in \texttt{ReachableDemandSeparation.lean}. This witness uses
complete evaluations of overlapping panels and establishes why global prefix
totals cannot serve as the cost state for that selective-demand protocol.
A policy retaining each query's membership, or the equivalent demand-profile
quotient, computes both states' bills exactly.

The distinction concerns information, not a prescribed data structure: a bitset or another lossless encoding of incidence is equally sufficient. A stronger prefix summary that preserves the needed conditional bills also avoids the counterexample. For query-selective evaluation, let $S$ denote the next query subset and $X$ the completed workflow. The prediction target is
\begin{equation}
 \begin{split}
 m_t(p,a)=\mathbb E\Bigl[
   &\sum_{q\in S}\sum_{u\preceq X,\,u\ne\varnothing}
       \ind\{(q,u)\notin K_t\}w_{q,u}\\[-2pt]
   &+h_t(S,X)\mid\mathcal H_t,p,a\Bigr].
 \end{split}
 \label{eq:joint-panel-cost}
\end{equation}
The joint law of $S$ and $X$ matters; separate query and path marginals need not determine their resource-inclusion probabilities. This identity follows by summing the conditional inclusion probability of each uncached resource. The full-panel estimator below sets $S=\mathcal Q$.

Executing a completion also changes later candidates' bills by publishing their shared resources. Equation~\eqref{eq:cacheforecast} quantifies this effect. A full next-evaluation estimator uses the current graph; a planner that deliberately invests in reusable work additionally needs remaining budget and a value for the resulting cache state. Neither an accurate next-bill estimate nor a logarithmic quality bound alone establishes such planning optimality.

Propositions~\ref{prop:graph-compression}--\ref{prop:graph-separation} identify when query incidence affects cost-aware allocation. The main experiment uses the fixed-full-panel compression case. Section~\ref{sec:finite-budget} connects cost errors to decisions and evaluation coverage. Theorem~\ref{thm:shared} extends the quality bound to overlapping fixed panels with separate reward backups and a shared cache; this is a theoretical extension beyond the main one-tree protocol.
